%=============================================================================
\section{Discussion}
\label{sec:discussion}
%=============================================================================

\subsection{Topological simplification as a dimension of IBD dysbiosis}

The central finding of this study---that IBD-associated microbiomes show
dramatically lower \Hone\ complexity than healthy controls---is consistent
with the known compositional features of IBD but reveals a structurally
distinct dimension of dysbiosis.
Compositional studies have documented the depletion of butyrate-producing taxa
(\textit{Faecalibacterium prausnitzii}, \textit{Roseburia spp.}) and the
expansion of Proteobacteria in IBD~\citep{petersen2017defining}.
These changes would be expected to reduce the density and connectivity of the
co-occurrence network, which is precisely what our topological metrics capture:
fewer loops (H1 count $-1.34$ SD lower), shorter loop lifetimes (mean lifetime
$-1.65$ SD), and a lower peak number of simultaneously active loops (max
Betti-1 $-2.07$ SD).

Importantly, topological reduction is \emph{not} the same as diversity
reduction.
Shannon entropy quantifies evenness of abundance; our loop metrics quantify
the multi-scale co-occurrence structure.
A community could have high Shannon diversity but few co-occurrence loops if
taxa co-occur independently; conversely, a low-diversity community with strong
pairwise dependencies could exhibit rich loop structure.
The relationship between topological and diversity metrics in this dataset
remains to be characterized formally (see Limitations).

The magnitude of effect sizes ($d = 0.75$--$2.38$) is notably large for a
microbiome study.
We attribute this partly to the paired bootstrap design: by repeatedly drawing
$n = 100$ from each group and comparing the resulting topological features
across 200 iterations, we are effectively estimating the typical topological
difference between a random IBD patient sample and a random healthy sample,
rather than comparing two fixed community-level representations.
The bootstrap variance captures within-group heterogeneity, which tends to
reduce effect sizes; that the effects remain large under this conservative
design strengthens confidence in the finding.

\subsection{Diet signal: structural reorganization rather than enrichment}

The diet comparison revealed an intriguing pattern: plant-based eaters had
\emph{fewer} H1 features than omnivores ($d = -0.34$) but \emph{longer-lived}
loops (mean lifetime $d = +0.65$; max lifetime $d = +0.32$).
This counterintuitive result---fewer but more persistent loops---suggests that
plant-based microbiomes are not simply richer in co-occurrence structure;
rather, they exhibit a different structural geometry.
A possible biological interpretation is that the higher fiber intake
characteristic of plant-based diets selectively enriches for a smaller number
of cross-feeding consortia with strong metabolic dependencies (high loop
persistence), while reducing the overall number of weaker, transient
co-occurrences (lower total count).

This interpretation is speculative and should be tested with mechanistic data.
The partial attenuation of the diet signal after matching (3/6 vs.\ 5/6
features surviving FDR) suggests that demographic factors---perhaps
correlated lifestyle variables---account for some of the difference.

\subsection{Antibiotic recency: insufficient evidence for a topological claim}

We find no robust topological signature of recent antibiotic use.
Only one feature (max Betti-1) survived FDR correction in full-group analyses,
and the matched analysis at $N = 80$ produced zero FDR-significant features.
We do not interpret this as evidence that antibiotics have no effect on
microbiome topology; rather, the \agp\ antibiotic field aggregates use
across a 6-month window and across all antibiotic classes, which almost
certainly dilutes any class- or dose-specific effect.
A more targeted analysis---comparing pre- and post-treatment samples, or
stratifying by antibiotic class---would be required to evaluate the topological
impact of antibiotics.

\subsection{Limitations}
\label{subsec:limitations}

\paragraph{Cross-sectional design.}
All \agp\ samples are single time-point observations.
We cannot infer whether topological simplification precedes \ibd\ onset,
accompanies active disease, or persists during remission.
Longitudinal designs with pre/post samples would be required to address
directionality.

\paragraph{16S rRNA resolution.}
The \agp\ BIOM file contains deblurred ASVs, which were not collapsed to
genus level in our analysis.
Species and strain-level interactions, which may contribute importantly to
loop structure, are not resolved.

\paragraph{No comparison with conventional diversity metrics.}
We did not compare topological metrics against Shannon diversity or Bray--Curtis
dissimilarity in the same samples.
Whether topological features provide information \emph{beyond} standard metrics
is an important question that we leave for follow-up work.

\paragraph{Single dataset.}
All results are from the \agp.
The \agp\ is a self-selected citizen-science cohort with potential
ascertainment bias; replication in an independent clinical cohort
(e.g., curatedMetagenomicData, iHMP) is needed before clinical claims
can be made.

\paragraph{No mechanistic interpretation.}
We do not identify which specific taxa drive the loop-structure differences.
Taxon-level attribution (e.g., identifying which cross-feeding pairs form the
persistent loops) would require further analysis.

\paragraph{Bootstrap non-independence.}
The 200 paired iterations are not fully independent because they draw
(without replacement) from the same finite pool of samples.
Effective sample size is therefore less than 200.
The sign-flip permutation test accounts for this to some degree (it permutes
within the same set of deltas), but formal correction for the effective
sample size was not applied.
